\documentclass[../../../main.tex]{subfiles}
\begin{document}

\subsection{Data Analysis}

Artifacts are signals that make recording more difficult and may hamper the analysis of the brain activity recorded with EEG. Such artifacts can be divided into two categories: neurophysiological artifacts and environmental noise \cite{Corsi2023}.

Neurophysiological artifacts correspond to noise generated by the subjects themselves, whether it is voluntary or not. The brain is far from being the only organ that generates electromagnetic activity. In particular, the eyes and heart produce electromagnetic activity. As a result, the main neurophysiological artifacts are related to cardiac activity and ocular activity (blinks and saccades), and can be visually seen during an EEG recording. A possible way to reduce ocular artifacts is to instruct the subject to avoid moving their eyes and, for short recordings only, to avoid eye blinking \cite{Corsi2023}.

Another neurophysiological artifact, can be induced by the subject's voluntary motion. Motion engenders muscular activity that can distort the recorded brain signals. Typical examples are jaw clenching and swallowing. A possible way to reduce the muscular artifacts is to instruct the subjects to remain as quiet as possible and to avoid moving their jaws \cite{Corsi2023}.

Environmental noise refers to artifacts generated by the environment that surrounds the experimental setup. They can be linked with mechanical vibrations like the presence of a tramway nearby, or associated with power lines occurring at 50 or 60 Hz. Use of a Faraday cage (shielded room) may help to achieve the absolute best signal-to-noise ratio. Modern EEG equipment, especially systems with active electrodes, can also reduce noise. Additionally, system noise refers to artifacts generated by the sensors themselves that could be broken \cite{Corsi2023}.

Preprocessing for artifact removal is probably the most crucial part when analyzing EEG data. The point here is to remove noise without eliminating information of interest of distorting the signal. The first thing to do is to extensively study the dataset, in particular, inspecting the EEG signals but also the associated broadband power spectra. This preliminary step enables identification of most artifacts and, most importantly, if they have a specific temporal and frequency signature (presence of periodic artifacts) \cite{Corsi2023}.

From this point, it is possible to choose a specific strategy to remove the observed noise. In the case of cardiac and ocular artifacts, given their clear pattern, an efficient way to isolate and reduce them consists of applying independent component analysis (ICA). In the case of artifacts at a specific frequency, like power like noise at 50 or 60 Hz, one can consider applying notch filters. In the case of muscular activity, applying a low pass filter with a cutoff frequency at 40 Hz can be of interest \cite{Corsi2023}.

% Signal artifacts
\begin{figure}[!htbp]
\centering
\includegraphics[width=0.6\linewidth]{appendix/methods/datasets/assets/eeg/signal-artifacts.png}
\caption[EEG signal artifacts]{Example of EEG signal artifacts \autocite{Corsi2023}.}
\label{Figure:SignalArtifacts}
\end{figure}

It is possible to directly analyze the signals recorded by the sensors. In such a case, one will say that the analysis is performed in the space of the sensors. However, it is possible to go one step further and estimate the activity within the brain. This processing step is called source reconstruction and consists of estimating neural correlations of signals locations. It can be performed when one wants to have access to a higher spatial resolution and to provide a more accurate description, and interpretation, of the neurophysiological activity occurring. For that purpose, both direct and inverse problems need to be solved \cite{Corsi2023}.

For the direct problem, we aim at modeling the electromagnetic field produced. For that purpose, it is necessary to consider both a physical model and a model that will predict how the electromagnetic field will be generated at the scalp. The simplest model is the spherical model, which considers the head as an ensemble of spheres (brain, cerebrospinal fluid, skull, scalp) each characterized by a given conductivity. More realistic models rely on geometrical reconstruction of the different layers that form the head tissues, directly extracted from magnetic resonance imaging (MRI). They consist of building meshes of the interfaces between different tissues \cite{Corsi2023}.

The purpose of the inverse problem is to reconstruct brain activity in the physical model. One of the main challenges is the nonuniqueness of its solution. In other words, a large number of brain activity patterns could generate the same signature detected at the sensor level. Therefore, some constraints or assumptions are essential to lead to a unique solution that best reflects the acquired data. Dipole modeling methods rely on modeling with a reduced number of equivalent dipoles and, therefore, are based on an a priori hypothesis of the required number of sources. Scanning methods consist of estimating the probability of presence of a current dipole inside each voxel \cite{Corsi2023}.
\end{document}
